% Section heading emitted by \introduction in main.tex
High-resolution satellite mapping of surface water is key for disaster management and coastal conservation \citep{Wu2020Ensemble}. High-resolution datasets like the Joint Research Centre (JRC) Global Surface Water product highlight rapid changes in terrestrial hydrology under climate change and human regulation \citep{Pekel2016Highresolution}. Recent advances leverage deep neural networks to segment surface water from active microwave or optical imagery \citep{Bazi2021Vision, Phan2020Land}. 

However, these empirical models operate as static spatial mappings. They bypass intermediate land surface states (e.g., storage or soil saturation) and ignore temporal catchment memory \citep{Nearing2020What}. In coastal deltas, this is highly problematic because upstream reservoir regulation and canal irrigation networks introduce non-linear anthropogenic lags that decouple local precipitation from local wetness, leading to spatial overfitting and temporal representation collapse.

Although previous research has successfully applied remote sensing classifiers and deep learning to map deltaic inundation, many existing frameworks treat geomorphological constraints and meteorological drivers as joint statistical features. They frequently fail to separate their individual predictive weights or identify the boundaries where topographic controls and climate triggers decouple. Consequently, previous models are insufficient because combining these features masks spatial extrapolation errors in new basins and leads to representation collapse during dry seasons when coarse climate data cannot resolve sub-grid channel networks. This study is scientifically necessary to isolate these controls under a coordinate-withheld, leakage-free design, establishing where meteorological inputs may fail to resolve fine topographic structures. This study provides a rigorous, spatially independent benchmark for separating geomorphological and meteorological controls on surface water dynamics in a highly regulated deltaic basin. By doing so, we position our work against traditional empirical classifiers by highlighting the physical limits of data-driven models under spatial and temporal scale mismatches.

Unlike previous delta flooding studies that treat geomorphological and meteorological variables as joint predictors, this work separates their roles using a leakage-free design, a point-level spatial split, and residual spatial autocorrelation testing. This allows us to quantify where static terrain explains spatial flood boundaries and where coarse climate forcing may fail to resolve sub-pixel hydrological structure.

To achieve these objectives, we present five primary scientific contributions, each directly supported by our experimental design and statistical audits. First, we establish a leakage-free flood prediction benchmark that excludes concurrent satellite observations, ensuring that the models rely solely on static terrain and antecedent climate features. Second, we execute a terrain-meteorological feature ablation study: Terrain features account for 87.81\% of the Random Forest Gini feature importance, representing model-specific predictive weights rather than direct physical causality. Concurrently, we find that Terrain-only and Combined models are statistically indistinguishable, with overlapping bootstrap confidence intervals. Third, we characterize temporal routing lag and irrigation decoupling, identifying a multi-month routing delay in coastal wetlands and an out-of-phase divergence in croplands. Fourth, we perform a spatial autocorrelation residual audit, showing that Moran's I on validation residuals is statistically indistinguishable from zero, which suggests that our disjoint coordinate split successfully prevented spatial leakage. Fifth, we diagnose sequence model representation collapse, showing that LSTMs and Transformers default to dry-class predictions when spatial coordinates are withheld, highlighting a resolution mismatch. These results indicate that high-resolution geomorphological constraints explain the spatial organization of inundation boundaries under the evaluated feature set, whereas coarse meteorological forcing primarily regulates temporal activation.
